% !TEX root = ../main.tex
\pagebreak[0]
\section{Remarks and further directions}\label{sec:remarks-and-further-directions}

\subsection{Relation to previous work}\label{subsec:relation-to-previous-work}

Most of the tools used here are established. Subdividing boxes of
parameters until a sufficient condition holds on each is the classical
branch and bound method. Exact arithmetic in a quadratic field, support
separation of projected vertices, and the centring argument for centrally
symmetric bodies are standard in this context; separation by a supporting
line also underlies the global theorem of
\citet[Section~4, Theorem~17]{SteiningerYurkevichNopertV2}. Bernstein
coefficients as bounds for the range of a multivariate polynomial on a box
go back at least to \citet{Garloff1986}.

Their combination, however, is new in this context. The parametrisation
turns every comparison into a polynomial inequality with coefficients in
$\mathbb Q(\sqrt5)$, so that no trigonometric function has to be bounded
and every step can be checked exactly. Bernstein coefficients then serve as
the single sign test of the whole proof, from the domain inequalities to
the gap polynomials, and the zoom lemma extends that same test to the
neighbourhoods of touch configurations. Together they make the entire
argument a finite list of exact checks, which a certificate records and an
independent program can repeat.

What the RID required, and what we supply, is a way to eliminate the
neighbourhoods of touch configurations, where no strict inequality can
hold and ordinary subdivision cannot finish.
\citet{SteiningerYurkevichNopertV2} proved the Noperthedron Nopert with a
global and a local theorem, and identified an exotic viewing direction
of the RID where their local theorem does not apply; they indicated that its neighbourhood could be treated with
polynomial inequalities, and reported further difficult regions
\citep[Section~9.1]{SteiningerYurkevichNopertV2}. Here the aligned
configurations, the square and pentagon configurations and two
non-aligned families of touch configurations are all treated by one
statement, \cref{lem:zoom-lemma}. The idea is to recognise sets of
configurations that contain no fit for a simple reason, change coordinates
so that some coordinates measure the distance from such a set, and divide
a necessary inequality exactly by a power of those distances before
testing its sign. This is reminiscent of blowing up a singular point in
algebraic geometry, although we use nothing from that theory. The
specific ingredients for the RID are the symmetry-reduced domain with its
fold inequalities, the observation that the non-aligned touch families lie
in planes where the plug reaches exactly as far as the hole in one screen
direction, and the explicit zoom covers of \cref{app:zoom-covers}.
\Cref{thm:main} settles the conjecture of \citet{SteiningerYurkevich2023}
that the RID is not Rupert.

\subsection{Extensions and possible automation}\label{subsec:extensions-and-possible-automation}

\paragraph{What transfers.}
The argument uses few properties of the RID beyond its exact vertex
coordinates, its central symmetry and its symmetry group. For another
centrally symmetric polyhedron with vertex coordinates in a real number
field, the Global component applies unchanged, the Local component applies
in the same form with new covers, and the Domain component needs that
polyhedron's own reduced domain. The Exotic component's covers are
specific to the RID, but their construction is not: each consists of a
no-fit set, an affine change to base and offset coordinates, and the boxes
and radii of a tube or point zoom. Much of the work is already automatic.
The zoom cells and their witnesses were found by a program, the parameters
of each cover amount to a few numbers, and exact arithmetic decides the
equalities on which the no-fit sets rest. The step that required insight
was finding the touch configurations and the no-fit sets containing them.
A natural next step is to apply the method to other polyhedra conjectured
to be Nopert.

\paragraph{Towards a decision procedure.}
We expect these techniques to lead to a practical algorithm that decides
Rupert's property for centrally symmetric polyhedra, once the special cases
can be enumerated algorithmically. The missing ingredient is a finite
procedure that finds, for a given polyhedron, every family of touch
configurations together with a no-fit set containing it, and the points on
these families where a second distance must be divided out. The aligned
configurations are always touch configurations, and the exotic viewing
directions of the Local component might be enumerated: if the component
fails only where faces are seen edge-on, the candidates are finitely many
intersections of great circles on the viewing sphere. Non-aligned touch
families, such as the arcs of \cref{sec:arc-configurations}, need a separate
analysis, and we do not know whether they always lie in no-fit sets as
simple as the arc planes. Given such an enumeration, the remaining steps
are mechanical: placing zoom covers around the families, searching for cells
and witnesses, and running the subdivision with the resulting collection,
which would also detect a fit if one existed
(\cref{subsec:robustly-detecting-potential-rupertness}). We expect this to
be practical for polyhedra whose vertex coordinates are algebraic of low
degree, whose symmetry group is large enough to make the reduced
configuration space small, and which are regular in other respects, as the
RID is; how far it extends is a question for further work.

The checks performed on the certificate are few and elementary, and
together with the lemmas of
\cref{sec:geometric-formulation,sec:symmetry-reduced-domain,sec:reduction-to-exact-arithmetic,sec:global-component,sec:local-component,sec:square-and-pentagon-configurations,sec:arc-configurations,sec:endpoint-and-crossing-configurations}
they would be natural candidates for formalisation in a proof assistant.

\subsection{History of this work}\label{subsec:history-of-this-work}

The author has worked on this problem since 2024, as a self-motivated
project pursued in free time, over roughly 1,500 hours. Already in 2024, the
author developed, independently, the two methods on which the proof for the
Noperthedron by \citet{SteiningerYurkevichNopertV2} rests: the global
separation theorem, which became the Global component, and a version of
their local theorem around aligned configurations, which became the Local
component. These first versions used an angular parametrisation with
transcendental functions, less practical than the one used here, together
with interval arithmetic. The author also knew at that stage about the
exotic configurations of the RID, where neither method applies, and
suspected that a polyhedron could be constructed for which the two
components alone would suffice. Nevertheless, the author chose to pursue
the case of the RID itself, whose exotic configurations these methods proved
hard to reach. In the meantime, \citet{SteiningerYurkevichNopertV2}
constructed the Noperthedron, the first polyhedron proved to be Nopert.

The next approach treated each exotic case with its own collection of
polynomial inequalities, designed around the geometric features of that
case. It was sufficient, but considerably more complicated: the proof was
far more involved, and generating the full certificate required thousands
of times as much computation as the proof presented here. The zoom lemma of
\cref{subsec:the-zoom-lemma} then unified all these
cases, simplifying both the theoretical foundations of the proof and its
implementation. It also contains the earlier components: the Local
component is its simplest instance with a distance divided out, and the
Domain and Global tests are the same idea with nothing divided out.

To complete the work without further delay, AI tools were used to produce
a well-tested implementation of the final pipeline, to simplify the
author's research pipelines, which had grown to more than 100,000 lines of
code over the years, and to help write this article. The ideas, the
mathematical arguments and the design of the pipeline are the author's.
